\begingroup
\renewcommand{\arraystretch}{1.35}
\InterpolationArticle{3}
Let $X$ be a function of the indeterminate $x$ of the form
\[
\alpha+\beta x+\gamma xx+\delta x^3+\text{etc.}
\]
which does not continue to infinity, but breaks off without going beyond the term containing $x^{m-1}$, so that the number of coefficients is no greater than $m$. Then, if for $m$ distinct values of $x$, namely $a,b,c,d\ldots$, the corresponding values of $X$ are known, namely $=A,B,C,D\ldots$, the value of $X$ corresponding to any other value of $x$ can be obtained neatly from these in the following way. Let $t$ be the new value of $x$, and $T$ the corresponding value of $X$, so that the following $m+1$ equations hold
\[
\begin{aligned}
A&=\alpha+\beta a+\gamma aa+\delta a^3+\ldots\\
B&=\alpha+\beta b+\gamma bb+\delta b^3+\ldots\\
C&=\alpha+\beta c+\gamma cc+\delta c^3+\ldots\\
D&=\alpha+\beta d+\gamma dd+\delta d^3+\ldots\\
&\text{etc.}\\
T&=\alpha+\beta t+\gamma tt+\delta t^3+\ldots
\end{aligned}
\]
\endgroup
\par
\begingroup
\renewcommand{\arraystretch}{1.45}
\SourceLeaf{282}
Let these equations be multiplied respectively by
\[
\begin{gathered}
\frac{1}{(a-b)(a-c)(a-d)\mathinner{.\,.\,.\,.}(a-t)}\\[4pt]
\frac{1}{(b-a)(b-c)(b-d)\mathinner{.\,.\,.\,.}(b-t)}\\[4pt]
\frac{1}{(c-a)(c-b)(c-d)\mathinner{.\,.\,.\,.}(c-t)}\\[4pt]
\frac{1}{(d-a)(d-b)(d-c)\mathinner{.\,.\,.\,.}(d-t)}\\
\text{etc.}\\
\frac{1}{(t-a)(t-b)(t-c)\mathinner{.\,.\,.\,.}}
\end{gathered}
\]
and let the products be added to give
\[
\begin{aligned}
&\frac{A}{(a-b)(a-c)(a-d)\mathinner{.\,.\,.}(a-t)}\\
&{}+\frac{B}{(b-a)(b-c)(b-d)\mathinner{.\,.\,.}(b-t)}\\
&{}+\frac{C}{(c-a)(c-b)(c-d)\mathinner{.\,.\,.}(c-t)}\\
&{}+\frac{D}{(d-a)(d-b)(d-c)\mathinner{.\,.\,.}(d-t)}\\
&{}+\text{etc.}\\
&{}+\frac{T}{(t-a)(t-b)(t-c)(t-d)\mathinner{.\,.\,.}}=W
\end{aligned}
\]
Then, from Article 1, where $m$ denoted what for us here is $m+1$, it is easy to conclude that $W=0$; thus multiplication by $(t-a)(t-b)(t-c)\mathinner{.\,.\,.}$ gives
\[
\begin{aligned}
T={}&\frac{(t-b)(t-c)(t-d)\mathinner{.\,.\,.}}{(a-b)(a-c)(a-d)\mathinner{.\,.\,.}}A\\
&{}+\frac{(t-a)(t-c)(t-d)\mathinner{.\,.\,.}}{(b-a)(b-c)(b-d)\mathinner{.\,.\,.}}B\\
&{}+\frac{(t-a)(t-b)(t-d)\mathinner{.\,.\,.}}{(c-a)(c-b)(c-d)\mathinner{.\,.\,.}}C\\
&{}+\frac{(t-a)(t-b)(t-c)\mathinner{.\,.\,.}}{(d-a)(d-b)(d-c)\mathinner{.\,.\,.}}D\\
&{}+\text{etc.}
\end{aligned}
\]
\InterpolationArticle{4}
The formula found in the preceding article has the property that, without any calculation, it is immediately evident that substituting for $t$ any one of the quantities $a,b,c,d\mathinner{.\,.}$ gives the
corresponding value $A,B,C,D\mathinner{.\,.}$. This is not its only recommendation: it is certainly by far the most convenient for practical use, at least whenever only a sin\SourceLeaf{283}%
gle value of $X$ is to be calculated from the given values. When several
are to be obtained, the following transformation of our formula may sometimes
be preferable
\[
\begin{aligned}
T={}&A+(t-a)\left(\frac{A}{a-b}+\frac{B}{b-a}\right)\\
&{}+(t-a)(t-b)\left(\frac{A}{(a-b)(a-c)}+\frac{B}{(b-a)(b-c)}+\frac{C}{(c-a)(c-b)}\right)\\
&{}+(t-a)(t-b)(t-c)\left(\frac{A}{(a-b)(a-c)(a-d)}+\frac{B}{(b-a)(b-c)(b-d)}\right.\\
&\qquad\left.{}+\frac{C}{(c-a)(c-b)(c-d)}+\frac{D}{(d-a)(d-b)(d-c)}\right)\\
&{}+\text{etc.}
\end{aligned}
\]
(It is derived very easily from the formula of Article 3 by supposing that the number of quantities
$a,b,c,d\mathinner{.\,.\,.}$ increases successively from the single quantity $a$ to the two $a,b$, then to the three $a,b,c$, etc.,
and subtracting each value of $T$ arising in this way from the next.)
Let the coefficients be denoted as follows:
\[
\begin{gathered}\frac{A}{a-b}+\frac{B}{(b-a)},\quad
\frac{A}{(a-b)(a-c)}+\frac{B}{(b-a)(b-c)}+\frac{C}{(c-a)(c-b)}\text{ etc.}\\
=A',\ A''\text{ etc.}\end{gathered}
\]
\Needspace{23\baselineskip}
further, let $B',B''$ etc. denote what results from $A',A''$ etc. when $a,b,c\mathinner{.\,.\,.}$
are replaced respectively by $b,c,d\mathinner{.\,.\,.}$, and $A,B,C\mathinner{.\,.\,.}$ respectively by $B,C,D\mathinner{.\,.\,.}$; similarly, let $C',C''$ etc. denote what results from $A',A''$ etc. when $a,b,c\mathinner{.\,.\,.}$ are replaced respectively by $c,d,e\mathinner{.\,.\,.}$, and $A,B,C\mathinner{.\,.\,.}$ respectively by $C,D,E\mathinner{.\,.\,.}$, and so on. Then we shall have
\[
\begin{array}{lll}
A'=\dfrac{A-B}{a-b},&B'=\dfrac{B-C}{b-c},&C'=\dfrac{C-D}{c-d}\ \mathinner{.\,.\,.}\\[5pt]
A''=\dfrac{A'-B'}{a-c},&B''=\dfrac{B'-C'}{b-d},&C''=\dfrac{C'-D'}{c-e}\ \mathinner{.\,.\,.}\\[5pt]
A'''=\dfrac{A''-B''}{a-d},&B'''=\dfrac{B''-C''}{b-e},&C'''=\dfrac{C''-D''}{c-f}\ \mathinner{.\,.\,.}\\
\text{etc.}&&
\end{array}
\]
and
\[\begin{aligned}
T&=A+A'(t-a)+A''(t-a)(t-b)\\&+A'''(t-a)(t-b)(t-c)+\mathinner{.\,.}
\end{aligned}\]
The number of quantities $A',B',C'$ etc. to be calculated here will be $m-1$, the number of
quantities $A'',B'',C''$ etc. will be $m-2$, and so on.
\InterpolationArticle{5}
The following modification of the formula of Article 3 can also be used
to advantage. Let $\mathfrak{X}$ be an arbitrary polynomial in the indeterminate $x$, but not beyond the
\SourceLeaf{284}
$m-1^{\mathrm{th}}$ degree; and let $\mathfrak{A},\mathfrak{B},\mathfrak{C}.\mathfrak{D}\mathinner{.\,.\,.\,.}\mathfrak{T}$ be its values for $x=a,b,c,d\mathinner{.\,.\,.}t$
respectively. Then, by Article 1,
\[
\begin{aligned}
0={}&\frac{\mathfrak{A}}{(a-b)(a-c)(a-d)\mathinner{.\,.\,.}(a-t)}\\
&{}+\frac{\mathfrak{B}}{(b-a)(b-c)(b-d)\mathinner{.\,.\,.}(b-t)}\\
&{}+\frac{\mathfrak{C}}{(c-a)(c-b)(c-d)\mathinner{.\,.\,.}(c-t)}\\
&{}+\frac{\mathfrak{D}}{(d-a)(d-b)(d-c)\mathinner{.\,.\,.}(d-t)}\\
&{}+\text{etc.}\\
&{}+\frac{\mathfrak{T}}{(t-a)(t-b)(t-c)(t-d)\mathinner{.\,.\,.}}
\end{aligned}
\]
From this it is easy to deduce
\[
\begin{aligned}
T={}&\mathfrak{T}\\
&{}+\frac{(t-b)(t-c)(t-d)\mathinner{.\,.\,.}}{(a-b)(a-c)(a-d)\mathinner{.\,.\,.}}(A-\mathfrak{A})\\
&{}+\frac{(t-a)(t-c)(t-d)\mathinner{.\,.\,.}}{(b-a)(b-c)(b-d)\mathinner{.\,.\,.}}(B-\mathfrak{B})\\
&{}+\frac{(t-a)(t-b)(t-d)\mathinner{.\,.\,.}}{(c-a)(c-b)(c-d)\mathinner{.\,.\,.}}(C-\mathfrak{C})\\
&{}+\frac{(t-a)(t-b)(t-c)\mathinner{.\,.\,.}}{(d-a)(d-b)(d-c)\mathinner{.\,.\,.}}(D-\mathfrak{D})\\
&{}+\text{etc.}
\end{aligned}
\]
\Needspace{12\baselineskip}
If the function $\mathfrak{X}$ is chosen so as to represent the values of the function $X$, namely $A,B,C,D\mathinner{.\,.\,.}$,
closely, we gain the advantage that the coefficients
$\frac{(t-b)(t-c)(t-d)\mathinner{.\,.\,.}}{(a-b)(a-c)(a-d)\mathinner{.\,.\,.}}$ etc.
need be multiplied only by small quantities. A constant quantity can also be taken for $\mathfrak{X}$,
for example one of $A,B,C,D\mathinner{.\,.\,.}$; in this case one term disappears from the value of
$T$. Or $\mathfrak{X}$ can be determined so that two or more of these quantities
are represented by $\mathfrak{X}$, in which case just as many terms disappear from $T$.
\InterpolationArticle{6}
What precedes rests on the assumption that the function $X$ does not go beyond the power
$x^{m-1}$, or that the $m^{\mathrm{th}}$ difference vanishes together with the higher differences; in
this case the method of interpolation is rigorously correct. If, however, that assumption
does not hold, interpolation aims only to find, in place of the function $X$, the \textit{simplest} function that satisfies the given values $A,B,C,D\mathinner{.\,.\,.}$. Now,
to enable us to judge the error arising from the neglected higher differences,
\SourceLeaf{285}
let the terms in \(X\) after the power \(x^{m-1}\) be \(\mu x^m+\nu x^{m+1}+\text{etc.}\). Thus, in Article 3,
the relation is not \(W=0\) but, as follows from Article 1, \(W=\mu+\nu(a+b+c+d+\mathinner{.\,.\,.}+t)+\text{etc.}\)
Hence the following must still be added to the value of \(T\) given there
\[
(t-a)(t-b)(t-c)(t-d)\mathinner{.\,.\,.}\times\{\mu+\nu(a+b+c+d+\mathinner{.\,.\,.}+t)+\text{etc.}\}
\]

\endgroup
